\section{A qualitative cutoff for tails of size five and six}
\label{sec:factorwise-five-six-tail}

\paragraph{Status and exact scope.}
HD-FT07--08, reserved question Q2dg, is a new derivation by
\texttt{/root/g34\_eightfolds}.  It uses the already checked character
polynomial of D-AFT01, the actual quotient interface of HD-FT01, and
the published arbitrary-fixed-genus theorem of Monderer--Neftin.
The cutoff below is qualitative: no numerical value of the published
constant \(N_g\) for the new, large values of \(g\) is asserted.
In particular it is not substituted into the effective FQ bound.
Complete author and separate root whole-proof readings pass. Root also
read the full source/status record, entire checker and first output;
the published arbitrary-genus passages were freshly reopened.
No numerical arbitrary-genus cutoff or new root execution is claimed.

Retain exactly (FT.1)--(FT.2): an actual connected finite-group cover
of the sphere, an actual quotient onto \(P=A_n\) or \(S_n\), an actual
factorization \(V=W\otimes U\) up to scalars, and genuine inertia with
\[
 e=\dim V,\qquad
 \sum_i\operatorname{rank}(1-T_i)=2e+6.
 \tag{FT56.1}
\]
The tensor partner and its field and degree are arbitrary.  No twelve-mark
condition is imposed on the projected permutations.  Let
\(\mu\vdash k\), where \(k\in\{5,6\}\), and put
\(W=S^{(n-k,\mu)}\), or its irreducible ordinary restriction to \(A_n\).
Transposition of the full partition, hence a sign twist in the symmetric
case, is allowed.  The bounds use absolute characters and projective
support, so this changes none of them.

Write the exact stable character polynomial as
\[
 q_\mu=C X_1^k+R,\qquad C=f^\mu/k!,\qquad
 L=\|R\|_1,\qquad M=1+L/C.
 \tag{FT56.2}
\]
Its weighted degree is at most \(k\), with weight \(i\) on \(X_i\);
\(R\) has ordinary degree at most \(k-1\).  These are the precise
premises of D-AFT01, from
\cite[Proposition~I.1 and complete proof, pp.~9--11]{HD:GarsiaGoupil2009}.
In particular \(M\ge1\).

\subsection*{HD-FT07. A cubic comparison with a cubic-order localized error}

\begin{lemma}[HD-FT07]
Suppose \(k=5\) or \(6\), and \(n\ge256M\).  Put
\[
 a=\dim W,\qquad
 \varepsilon=\frac{2M(M+1)}{n^3}.
 \tag{FT56.3}
\]
Then \(a\ge2n^3\), and for every permutation \(x\), with
\(f=X_1(x)\), one has
\[
 \frac{|\chi_W(x)|}{a}\le
 \begin{cases}
 (f/n)^3,&f\ge n/2,\\
 (f/n)^3+\varepsilon,&f<n/2.
 \end{cases}
 \tag{FT56.4}
\]
This includes every power of every inertia and every target eigenvalue
in the Fourier consequence used below.
\end{lemma}

\begin{proof}
Set \(t=f/n\) and \(m=n-f\).  The ordinary degree of \(R\) gives
\[
 a\ge Cn^k-Ln^{k-1}\ge Cn^k(1-M/n)\ge\frac{n^k}{2k!}.
 \tag{FT56.5}
\]
Since \(n\ge256\), \(n^{k-3}\ge4k!\) for both \(k=5,6\),
which proves \(a\ge2n^3\).  Also \(n>2k\), as required for the
stable partition and its non-self-conjugate alternating restriction.

For \(f\ge n/2\), the argument of HD-FT04 applies with these same
coefficients, and we repeat its inequalities to specify the threshold.
The character is positive, since
\(Cf^k-Ln^{k-1}>0\) follows from \(n>2^kL/C\).
For the remainder,
\[
 |R(f,X_2,\ldots)-R(n,0,\ldots)|
       \le kLm n^{k-2}.
 \tag{FT56.6}
\]
A pure \(f^j\) term satisfies the derivative estimate, and a mixed
term has a non-\(f\) cycle count at most \(m\); these observations
prove (FT56.6) term by term.  Furthermore \(1-t^3\le3m/n\).
Consequently
\[
 q_\mu-a t^3
 \le -Cn^k(t^3-t^k)+(k+3)Lm n^{k-2}\le0.
 \tag{FT56.7}
\]
Indeed \(t^3-t^k\ge(1-t)/8\) for \(1/2\le t\le1\) and \(k\ge4\),
and \(n\ge256M\ge8(k+3)L/C\).  This proves the first case of
(FT56.4), including \(f=n\).

Suppose \(f<n/2\).  For a monomial of \(R\), write \(j\) for the
exponent of \(f\) and \(b\) for the sum of the other exponents.
Then \(j+2b\le k\) and \(j+b\le k-1\).
After division by \(n^k\), its absolute value is at most
\(t^3/n\) if \(j\ge3\), at most \(t^2/n^2\) if \(j=2\),
at most \(t/n^2\) if \(j=1\), and at most \(n^{-3}\) if \(j=0\).
Thus the following deliberately coarse simultaneous bound holds:
\[
 \frac{|R(f,X_2,\ldots)|}{Cn^k}
 \le M\left(\frac{t^3}{n}+\frac{t^2}{n^2}
                   +\frac{t}{n^2}+\frac1{n^3}\right).
 \tag{FT56.8}
\]
For \(z\ge0\) and \(M\ge1\), the elementary arithmetic--geometric
mean inequality gives
\[
 3Mz\le z^3+2M^{3/2}\le z^3+2M^2.
 \tag{FT56.9}
\]
Apply this with \(z=nt\).  Because \(t^2\le t/2\) and
\(t^k\le t^3/4\) on \(0\le t<1/2\), (FT56.8)--(FT56.9) imply
\[
 \frac{|q_\mu|}{Cn^k}
 \le\left(\frac34+\frac Mn\right)t^3
                  +\frac{M(M+1)}{n^3}.
 \tag{FT56.10}
\]
Here the two linear/quadratic terms contribute at most
\(t^3/2+M^2/n^3\).  Divide by the lower bound
\(a/(Cn^k)\ge1-M/n\).
Since \(M/n\le1/256\le1/8\), the cubic coefficient is at most one,
and \((1-M/n)^{-1}\le2\).  This proves the second case of (FT56.4).
\end{proof}

\subsection*{HD-FT08. Arbitrary partners and a qualitative natural-degree cutoff}

For each \(\mu\) as above choose the integer
\[
 G_\mu=\left\lceil2M(M+1)+2\right\rceil.
 \tag{FT56.11}
\]
Let \(N_g\) be the constant in the published
\cite[Theorem~1.2]{HD:MondererNeftin2022}, at the fixed genus bound
\(g=G_\mu\), and set
\[
 N_\mu^{(56)}
   =\max\{\lceil256M\rceil,N_{G_\mu},256\}.
 \tag{FT56.12}
\]
Only existence of \(N_{G_\mu}\) is used; the known explicit constant
for genus one is not asserted to work at \(G_\mu\).

\begin{theorem}[HD-FT08]
No actual source (FT56.1) with the indicated factor \(W\) exists when
\(n>N_\mu^{(56)}\).  This holds with arbitrary scalar factors,
arbitrary coefficient fields and arbitrary factorwise tensor partners.
There are only eighteen partitions of sizes five or six, so taking
the maximum of these eighteen constants gives a finite qualitative
cutoff for both tails.
\end{theorem}

\begin{proof}
Call an inertia image \(x_i\) bad if it fixes fewer than \(n/2\)
natural letters.  If \(x_i\) is not bad, every power fixes at least
the same letters, so the first case of (FT56.4) applies to every power.
Cyclic Fourier inversion therefore gives, for every image,
\[
 \frac{\operatorname{rank}(1-T_i)}e
 \ge\frac{\nu(W(x_i))}{a}
 \ge\frac{\operatorname{ind}(x_i,\{1,\ldots,n\}^3)}{n^3}
                    -\varepsilon\,{\bf1}_{\{i\ {\rm bad}\}}.
 \tag{FT56.13}
\]
The first inequality is tensor spectral convolution; it handles
all target scalars and does not require reality of a chosen inertia.

There are at most ten bad entries.  To see this without assigning
any mark count to \(W\), use \(u^3\le u\) for \(0\le u\le1\)
in the same Fourier average.  For a bad image,
\[
 \frac{\nu(W(x_i))}{a}
 \ge\frac{\operatorname{ind}(x_i)}n-\varepsilon
 >\frac14-\varepsilon>\frac15.
 \tag{FT56.14}
\]
Indeed \(\operatorname{ind}(x_i)\ge(n-f(x_i))/2>n/4\), and
\[
 \varepsilon\le\frac{4M^2}{n^3}
       \le\frac1{16384n}<\frac1{20}.
 \tag{FT56.15}
\]
On the other hand \(e\ge a\ge2n^3\) implies
\(2+6/e\le2+3/n^3<11/5\).
Eleven bad entries would strictly violate this actual total budget.

Consider the actual quotient curve associated with the stabilizer
of an unordered three-set in \(P\).  Its stabilizer has order at
most \(6(n-3)!\); it cannot contain \(A_{n-2}\), of order
\((n-2)!/2\), because \(n\ge256>14\).
It also cannot contain \(A_{n-1}\).
Thus Monderer--Neftin's Theorem~1.2 implies that its genus
\(\widehat\gamma_3\) exceeds \(G_\mu\) when \(n>N_{G_\mu}\).
This application uses the actual \(P\)-cover and its degree-\(n\)
natural quotient, not a proposed numerical branch profile.
The ordered-distinct-triple quotient has degree six over this
unordered-triple curve.  Consequently its genus satisfies
\[
 \gamma_3\ge6\widehat\gamma_3-5\ge6G_\mu+1.
 \tag{FT56.16}
\]

As in HD-FT01, the action on ordered triples with repetitions
has five orbits, with quotient genera
\(\gamma_1,\gamma_2,\gamma_2,\gamma_2,\gamma_3\).
Riemann--Hurwitz on these actual connected covers gives
\[
 \sum_i\operatorname{ind}(x_i,\{1,\ldots,n\}^3)
 =2(n^3-5)+2(\gamma_1+3\gamma_2+\gamma_3)
 \ge2n^3+12G_\mu-8.
 \tag{FT56.17}
\]
Summing (FT56.13), using the at-most-ten bound and (FT56.11), yields
\[
 \begin{split}
 2+\frac6e
 &\ge 2+\frac{12G_\mu-8-20M(M+1)}{n^3}\\
 &\ge 2+\frac{4M(M+1)+16}{n^3}
  >2+\frac3{n^3}\ge2+\frac6e,
 \end{split}
 \tag{FT56.18}
\]
a contradiction.  This proves the theorem.
\end{proof}

\paragraph{Remaining gate and source precision.}
Combined with HD-FT01--06, this removes unbounded tails of sizes
five and six from that actual factorwise source setup.  The only
unbounded ordinary extendible alternating labels left by these
particular arguments have tail size at most four.
This is a qualitative reduction: the finitely many possible intervals
between the existing effective cutoff and \(N_\mu^{(56)}\) have
not been enumerated, and no effective replacement for those constants
has been claimed.  No conclusion about a general complex coefficient
factor without the actual quotient (FT.2) is inferred.
The published arbitrary-genus theorem, not an unrefereed numerical
genus-growth estimate, is the new source input.
